% GENERATED FILE. Edit canonical lessons, then run npm run generate:books.
% canonical-source-hash: fnv1a64:9d88c5b7a2cabcb5
% canonical-lessons: IT-C350-siccome, IT-C350-poiche, IT-C350-percio, IT-C350-pertanto, IT-C350-dovuto

\chapter{Siccome e poiché}
\label{ch:it-siccome-e-poiche}
\hlchaptermodality{350}

\begin{chapteropening}
\textbf{By the end of this chapter:} \emph{I can link a cause to its result: since, as, so and therefore, and say what something is due to.}

Complete the last lesson of chapter 350: I can link a cause to its result: since, as, so and therefore, and say what something is due to.
\end{chapteropening}

\section[siccome]{siccome — since, as (at the start of a sentence)}

\label{lesson:IT-C350-siccome}



\begin{quote}
Before the new one: say the Italian for the factor, then the Italian for the reason.
\end{quote}

\subsection*{You'll want to know: siccome}
\textbf{siccome} — \textquotedblleft{}since, as (at the start of a sentence)\textquotedblright{}.

\begin{quote}
Siccome piove, prendiamo la macchina.
\end{quote}

\begin{quote}
\emph{Since it is raining, we will take the car.}
\end{quote}

\emph{Siccome} opens the sentence; \emph{perché} comes after the main clause.

\begin{cousinweb}
\emph{siccome} is built from \emph{sì} \textquotedblleft{}so\textquotedblright{}, inherited from Latin \emph{sīc} \textquotedblleft{}thus\textquotedblright{}, and \emph{come} \textquotedblleft{}as\textquotedblright{}, inherited from Latin \emph{quōmodo} \textquotedblleft{}in what way\textquotedblright{} (Lewis \& Short).
\end{cousinweb}

\begin{grammarlens}[title={What you've built}]
One more word for telling, explaining and describing at B1.
\end{grammarlens}

\subsection*{Your turn}
\begin{itemize}
\raggedright
  \item \emph{Say it:} \emph{siccome}
  \item \emph{Say it:} \emph{siccome}, once more
  \item \emph{Say it:} a sentence of your own with \emph{siccome}
  \item \emph{Recall:} say the Italian for the factor, then the Italian for the reason, then say \emph{siccome} again
\end{itemize}

\begin{tcolorbox}[breakable,colback=teal!4,colframe=teal!35!black,title={Before you move on}]
What does siccome mean? (\textquotedblleft{}Since, as (at the start of a sentence)\textquotedblright{}.) Say it once more.
\end{tcolorbox}
\section[poiché]{poiché — given that, since (formal)}

\label{lesson:IT-C350-poiche}



\begin{quote}
Before the new one: say the Italian for the reason, then the Italian for since.
\end{quote}

\subsection*{You'll want to know: poiché}
\textbf{poiché} — \textquotedblleft{}given that, since (formal)\textquotedblright{}.

\begin{quote}
Il museo è chiuso, poiché oggi è lunedì.
\end{quote}

\begin{quote}
\emph{The museum is closed, since today is Monday.}
\end{quote}

Written or formal; in speech \emph{perché} or \emph{siccome} are usual.

\begin{cousinweb}
\emph{poiché} is built from \emph{poi} \textquotedblleft{}then, after\textquotedblright{}, inherited from Latin \emph{post} \textquotedblleft{}after\textquotedblright{} (Lewis \& Short), and \emph{che} \textquotedblleft{}that\textquotedblright{}: first \textquotedblleft{}after that\textquotedblright{}, then \textquotedblleft{}since\textquotedblright{}.
\end{cousinweb}

\begin{grammarlens}[title={What you've built}]
One more word for telling, explaining and describing at B1.
\end{grammarlens}

\subsection*{Your turn}
\begin{itemize}
\raggedright
  \item \emph{Say it:} \emph{poiché}
  \item \emph{Say it:} \emph{poiché}, once more
  \item \emph{Say it:} a sentence of your own with \emph{poiché}
  \item \emph{Recall:} say the Italian for the reason, then the Italian for since, then say \emph{poiché} again
\end{itemize}

\begin{tcolorbox}[breakable,colback=teal!4,colframe=teal!35!black,title={Before you move on}]
What does poiché mean? (\textquotedblleft{}Given that, since (formal)\textquotedblright{}.) Say it once more.
\end{tcolorbox}
\section[perciò]{perciò — so, therefore}

\label{lesson:IT-C350-percio}



\begin{quote}
Before the new one: say the Italian for since, then the Italian for given that.
\end{quote}

\subsection*{You'll want to know: perciò}
\textbf{perciò} — \textquotedblleft{}so, therefore\textquotedblright{}.

\begin{quote}
Ho perso il treno, perciò sono arrivato tardi.
\end{quote}

\begin{quote}
\emph{I missed the train, so I arrived late.}
\end{quote}

A little more formal than \emph{quindi}; very frequent in writing and speech.

\begin{cousinweb}
\emph{perciò} is built from \emph{per} \textquotedblleft{}for\textquotedblright{} and \emph{ciò} \textquotedblleft{}this\textquotedblright{}, which comes from Latin \emph{ecce hoc} \textquotedblleft{}look, this\textquotedblright{} (Lewis \& Short): \textquotedblleft{}for this\textquotedblright{}, and so \textquotedblleft{}therefore\textquotedblright{}.
\end{cousinweb}

\begin{grammarlens}[title={What you've built}]
One more word for telling, explaining and describing at B1.
\end{grammarlens}

\subsection*{Your turn}
\begin{itemize}
\raggedright
  \item \emph{Say it:} \emph{perciò}
  \item \emph{Say it:} \emph{perciò}, once more
  \item \emph{Say it:} a sentence of your own with \emph{perciò}
  \item \emph{Recall:} say the Italian for since, then the Italian for given that, then say \emph{perciò} again
\end{itemize}

\begin{tcolorbox}[breakable,colback=teal!4,colframe=teal!35!black,title={Before you move on}]
What does perciò mean? (\textquotedblleft{}So, therefore\textquotedblright{}.) Say it once more.
\end{tcolorbox}
\section[pertanto]{pertanto — therefore, consequently (formal)}

\label{lesson:IT-C350-pertanto}



\begin{quote}
Before the new one: say the Italian for given that, then the Italian for so.
\end{quote}

\subsection*{You'll want to know: pertanto}
\textbf{pertanto} — \textquotedblleft{}therefore, consequently (formal)\textquotedblright{}.

\begin{quote}
Il negozio è chiuso ad agosto; pertanto non possiamo cambiare il regalo.
\end{quote}

\begin{quote}
\emph{The shop is closed in August; consequently we cannot exchange the present.}
\end{quote}

Formal and written. Careful: in some regions and older texts it can mean however.

\begin{cousinweb}
\emph{pertanto} is built from \emph{per} \textquotedblleft{}for\textquotedblright{} and \emph{tanto} \textquotedblleft{}so much\textquotedblright{}, inherited from Latin \emph{tantum} \textquotedblleft{}so much\textquotedblright{} (Lewis \& Short): \textquotedblleft{}for so much\textquotedblright{}, so \textquotedblleft{}therefore\textquotedblright{}.
\end{cousinweb}

\begin{grammarlens}[title={What you've built}]
One more word for telling, explaining and describing at B1.
\end{grammarlens}

\subsection*{Your turn}
\begin{itemize}
\raggedright
  \item \emph{Say it:} \emph{pertanto}
  \item \emph{Say it:} \emph{pertanto}, once more
  \item \emph{Say it:} a sentence of your own with \emph{pertanto}
  \item \emph{Recall:} say the Italian for given that, then the Italian for so, then say \emph{pertanto} again
\end{itemize}

\begin{tcolorbox}[breakable,colback=teal!4,colframe=teal!35!black,title={Before you move on}]
What does pertanto mean? (\textquotedblleft{}Therefore, consequently (formal)\textquotedblright{}.) Say it once more.
\end{tcolorbox}
\section[dovuto]{dovuto — due (to), caused (by); owed}

\label{lesson:IT-C350-dovuto}



\begin{quote}
Before the new one: say the Italian for so, then the Italian for therefore.
\end{quote}

\subsection*{You'll want to know: dovuto}
\textbf{dovuto} — \textquotedblleft{}due (to), caused (by); owed\textquotedblright{}.

\begin{quote}
Il ritardo è dovuto al traffico.
\end{quote}

\begin{quote}
\emph{The delay is due to the traffic.}
\end{quote}

\emph{Essere dovuto a} — to be due to. It is also the past participle of \emph{dovere}: \emph{ho dovuto aspettare}.

\begin{cousinweb}
\emph{dovuto} is the past participle of \emph{dovere}, inherited from Latin \emph{dēbēre} \textquotedblleft{}to owe, to be bound to\textquotedblright{} (Lewis \& Short): what is owed to a cause.
\end{cousinweb}

\begin{grammarlens}[title={What you've built}]
One more word for telling, explaining and describing at B1.
\end{grammarlens}

\subsection*{Your turn}
\begin{itemize}
\raggedright
  \item \emph{Say it:} \emph{dovuto}
  \item \emph{Say it:} \emph{dovuto}, once more
  \item \emph{Say it:} a sentence of your own with \emph{dovuto}
  \item \emph{Recall:} say the Italian for so, then the Italian for therefore, then say \emph{dovuto} again
\end{itemize}

\begin{tcolorbox}[breakable,colback=teal!4,colframe=teal!35!black,title={Before you move on}]
What does dovuto mean? (\textquotedblleft{}Due (to), caused (by); owed\textquotedblright{}.) Say it once more.
\end{tcolorbox}
